\chapter{Stochastic Calculus and the Diffusion Channel}
\label{app:stochastic}

The It\^o calculus, the Fokker--Planck equation, sampling by dynamics, and the
AR(1) chain in full, supporting Chapter~\ref{ch:background}.

\section{The AR(1) chain: the running example of the thesis}\label{sec:ar1}

We now introduce the specific Markov chain that the research chapters
solve completely: the first-order autoregressive (AR(1)) process, the
canonical Gaussian Markov chain \citep{brockwell1991time}. Let
$\alpha \in (-1, 1)$ and let $(\eta_k)$ be i.i.d.\
$\Nd(0, \sigma_\eta^2)$ innovations; define
\begin{equation}
  a_{k+1} \;=\; \alpha\, a_k + \eta_k ,
  \qquad k = 0, 1, \dots, K-2 .
  \label{eq:ar1}
\end{equation}
The transition kernel is Gaussian,
$M(a' | a) = \Nd(a'; \alpha a, \sigma_\eta^2)$, and the joint law of the
trajectory follows from \eqref{eq:chain-factorisation}. Throughout the
thesis we use the stationary normalisation
$\sigma_\eta^2 = 1 - \alpha^2$ and start the chain in its stationary law,
$a_0 \sim \Nd(0,1)$, so that every frame has unit marginal variance. A
trajectory $a = (a_0, \dots, a_{K-1})$ is \emph{one} data point, the
prototype of a video or a time series; its frame index $k$ is an internal
time, to be kept distinct from the diffusion time $t$ introduced in
Section~\ref{sec:ou}. Everything about this chain is computable in closed
form; the next toolbox derives the three facts used throughout.

\begin{toolbox}[tb:ar1]{Stationary distribution and autocovariance of AR(1)}
\emph{(i) Stationary variance.} Suppose $a_k \sim \Nd(0, \sigma^2)$ and
impose that $a_{k+1}$ has the same variance. Since $\eta_k$ is
independent of $a_k$,
\begin{equation*}
  \operatorname{Var}(a_{k+1})
  = \alpha^2 \operatorname{Var}(a_k) + \sigma_\eta^2
  \;\stackrel{!}{=}\; \operatorname{Var}(a_k)
  \quad\Longrightarrow\quad
  \sigma_\infty^2 = \frac{\sigma_\eta^2}{1 - \alpha^2}.
\end{equation*}
The fixed point exists precisely because $|\alpha| < 1$; the map
$\sigma^2 \mapsto \alpha^2 \sigma^2 + \sigma_\eta^2$ is a contraction, so
\emph{any} initial variance converges to $\sigma_\infty^2$ geometrically.
Since Gaussianity is preserved by the affine recursion, the stationary
law is $\pi = \Nd(0, \sigma_\infty^2)$, which equals $\Nd(0,1)$ under the
stationary normalisation.

\emph{(ii) Autocovariance.} In the stationary regime, for $d \geq 1$,
unroll the recursion $d$ times:
$a_{k+d} = \alpha^d a_k + \sum_{j=0}^{d-1} \alpha^{j} \eta_{k+d-1-j}$.
All innovation terms are independent of $a_k$, hence
\begin{equation*}
  \operatorname{Cov}(a_k, a_{k+d})
  = \alpha^d\, \operatorname{Var}(a_k)
  = \alpha^d\, \sigma_\infty^2 ,
\end{equation*}
and by symmetry
$\operatorname{Cov}(a_i, a_j) = \alpha^{|i-j|} \sigma_\infty^2$ for all
$i, j$: correlations decay \emph{geometrically} with the lag, with rate
$\alpha$. The covariance matrix of $K$ consecutive frames,
\begin{equation*}
  (\Sigz)_{ij} = \sigma_\infty^2\, \alpha^{|i-j|},
\end{equation*}
is symmetric Toeplitz, as stationarity requires
(Section~\ref{sec:stationarity}).

\emph{(iii) Detailed balance.} With $\pi = \Nd(0, \sigma_\infty^2)$ and
$M(a'|a) = \Nd(a'; \alpha a, \sigma_\eta^2)$, both sides of
\eqref{eq:detailed-balance} are proportional to
\begin{equation*}
  \exp\!\left[
    -\frac{a^2}{2\sigma_\infty^2}
    -\frac{(a' - \alpha a)^2}{2\sigma_\eta^2}
  \right]
  = \exp\!\left[
    -\frac{a^2 + (a')^2 - 2\alpha\, a a'}{2\sigma_\eta^2}
  \right],
\end{equation*}
where the exponent on the right is \emph{symmetric} in $(a, a')$. Hence
the stationary AR(1) chain is reversible: statistically, the movie of the
chain looks the same played in either direction. \hfill$\square$
\end{toolbox}

The AR(1) chain is thus a working miniature of everything this chapter
discusses: a Markov kernel with an explicit stationary law, geometric
convergence, reversibility, and, decisive for us, exact Gaussian algebra
at every step.

\section{Sampling by dynamics: MCMC and Langevin}\label{sec:mcmc}

Sampling inverts the logic of the previous sections. There, a kernel $M$
was given and we asked for its stationary law $\pi$. In sampling problems
we are given $\pi$, typically a Boltzmann distribution
$\pi(x) \propto e^{-\beta H(x)}$ known only up to its partition function,
and we must \emph{design} a kernel whose stationary law is $\pi$; the
ergodic theorem \eqref{eq:ergodic} then converts one long trajectory into
expectation values. This is how the negative phase of Boltzmann-machine
learning \eqref{eq:bm-learning} is estimated in practice, and the idea
originates, fittingly, in computational statistical mechanics
\citep{metropolis1953equation}.

The \emph{Metropolis--Hastings} chain \citep{metropolis1953equation,
hastings1970monte} moves from $x$ by proposing $x' \sim g(\cdot|x)$ from
any tractable proposal kernel and accepting the move with probability
\begin{equation}
  A(x \to x') \;=\; \min\!\left\{ 1,\;
    \frac{\pi(x')\, g(x \,|\, x')}{\pi(x)\, g(x' \,|\, x)}
  \right\},
  \label{eq:mh-acceptance}
\end{equation}
staying at $x$ otherwise. The target enters only through the \emph{ratio}
$\pi(x')/\pi(x)$, so the intractable partition function cancels:
Metropolis--Hastings samples from $e^{-\beta H}/Z$ while only ever
evaluating energy differences. The acceptance rule is engineered so that
the resulting kernel satisfies detailed balance
\eqref{eq:detailed-balance} with respect to $\pi$; the verification is a
short computation \citep[Chapter~7]{robert2004monte}. When $x$ is
high-dimensional but each conditional $\pi(x_i | x_{\setminus i})$ is
tractable, the \emph{Gibbs sampler} \citep{geman1984stochastic} cycles
through the coordinates resampling each from its conditional, a
Metropolis--Hastings move whose acceptance probability is identically
one. Gibbs sampling thrives exactly when the model has sparse conditional
structure: for a Markov chain, $\pi(x_k | x_{\setminus k})$ depends only
on the two neighbours, and for the bipartite RBM of
Section~\ref{sec:boltzmann-machines} each layer updates in one parallel
step. Inference cost follows graph structure, a principle that
Chapter~\ref{ch:graphs} makes systematic.

For continuous $x$ and differentiable energy, a different design
principle uses gradients. Consider the discrete-time update
\begin{equation}
  x_{k+1} \;=\; x_k + \epsilon\, \nabla \log \pi(x_k)
  \;+\; \sqrt{2\epsilon}\; \xi_k ,
  \qquad \xi_k \sim \Nd(0, I),
  \label{eq:ula}
\end{equation}
the (unadjusted) \emph{Langevin algorithm}. The drift pushes towards high
probability; the noise maintains exploration. Two observations prepare
what follows. First, \eqref{eq:ula} requires only the \emph{score}
$\nabla \log \pi$; the normalising constant is never needed. This is
precisely why the score is the natural object for generative modelling
(Chapter~\ref{ch:diffusion}), and why so much of this thesis is devoted
to computing scores exactly. Second, as $\epsilon \to 0$ the recursion
\eqref{eq:ula} becomes a stochastic differential equation, the Langevin
diffusion, whose stationary law is exactly $\pi$. Making that limit
precise requires the It\^o calculus and Fokker--Planck theory that occupy
the rest of the chapter, where the claim is verified by direct
computation.

\section{From ODEs to SDEs: the It\^o calculus}\label{sec:ode-to-sde}

An ordinary differential equation $\dot x = f(x, t)$ moves a point along
a deterministic velocity field. Physical systems, however, feel rapid,
unresolved forces (collisions, thermal agitation) whose cumulative effect
is well modelled by adding noise:
\begin{equation}
  \dd X_t \;=\; \underbrace{f(X_t, t)\,\dd t}_{\text{drift}}
  \;+\; \underbrace{g(t)\,\dd W_t}_{\text{diffusion}} .
  \label{eq:generic-sde}
\end{equation}
Here $W_t$ is a \emph{Wiener process} (standard Brownian motion): the
unique process with $W_0 = 0$, independent increments,
$W_t - W_s \sim \Nd(0, t-s)$, and continuous paths. Equation
\eqref{eq:generic-sde} is shorthand for the integral equation
$X_t = X_0 + \int_0^t f(X_s, s)\,\dd s + \int_0^t g(s)\,\dd W_s$, and the
whole subtlety lives in the last term: Brownian paths are nowhere
differentiable, so $\int g\,\dd W$ cannot be a Riemann--Stieltjes
integral. It\^o's construction \citep{ito1944stochastic} defines it as
the limit of left-endpoint sums,
\begin{equation}
  \int_0^t g(s)\, \dd W_s
  \;=\; \lim_{n \to \infty} \sum_{j=0}^{n-1}
        g(s_j)\,\big( W_{s_{j+1}} - W_{s_j} \big),
  \label{eq:ito-integral}
\end{equation}
where the evaluation at the \emph{left} endpoint $s_j$ makes the
integrand independent of the coming increment. Two consequences follow
immediately: the It\^o integral is a martingale (zero mean), and its
variance is given by the \emph{It\^o isometry},
$\E\big[ (\int_0^t g\,\dd W)^2 \big] = \int_0^t g(s)^2\,\dd s$, both
because cross terms $\E[g(s_j)(W_{s_{j+1}}-W_{s_j})]$ vanish by
independence.

It is worth recording why no convention-free definition is available, since
the same two sums reappear in Chapter~\ref{ch:ring}. Fix $[0,t]$ and a
partition of mesh tending to zero. The \emph{quadratic variation}
$Q_n = \sum_j (W_{s_{j+1}} - W_{s_j})^2$ has mean
$\sum_j \Delta s_j = t$ and variance $2\sum_j (\Delta s_j)^2$, which vanishes
in the limit, so $Q_n \to t$ in mean square. The \emph{total variation}
$V_n = \sum_j |W_{s_{j+1}} - W_{s_j}|$ instead diverges, because
$Q_n \le (\max_j |W_{s_{j+1}} - W_{s_j}|)\, V_n$ and the maximum increment
tends to zero by continuity: a bounded $V_n$ would force $Q_n \to 0$,
contradicting $Q_n \to t > 0$. A Riemann--Stieltjes integral $\int f \dd g$
requires $g$ to have bounded variation, precisely so that the limit does not
depend on where inside each subinterval $f$ is evaluated. For Brownian motion
that guarantee is absent, different evaluation points give genuinely different
limits, and a convention must be fixed.

The price of the construction is a new chain rule. Its origin is the
quadratic variation just computed: increments scale as
$\dd W_t \sim \sqrt{\dd t}$, so second-order terms in $\dd W$ are
first-order in $\dd t$ and cannot be discarded.

\begin{toolbox}[tb:ito-lemma]{It\^o's lemma, and where the extra term comes from}
Let $Y_t = \phi(X_t, t)$ with $\phi$ twice differentiable and $X_t$
solving \eqref{eq:generic-sde}. Taylor-expand $\phi$ to second order:
\begin{equation*}
  \dd Y
  = \frac{\partial \phi}{\partial t}\,\dd t
  + \frac{\partial \phi}{\partial x}\,\dd X
  + \frac{1}{2} \frac{\partial^2 \phi}{\partial x^2}\,(\dd X)^2
  + \dots
\end{equation*}
Substitute $\dd X = f\,\dd t + g\,\dd W$ and expand the square:
\begin{equation*}
  (\dd X)^2 = f^2 (\dd t)^2 + 2 f g\, \dd t\, \dd W + g^2 (\dd W)^2 .
\end{equation*}
Now use the multiplication table of It\^o calculus, which is nothing but
the scaling $\dd W \sim \sqrt{\dd t}$ made systematic:
\begin{equation*}
  (\dd t)^2 = 0, \qquad \dd t\, \dd W = 0, \qquad (\dd W)^2 = \dd t .
\end{equation*}
The last identity is the rigorous statement that the quadratic variation
of $W$ over $[0,t]$ equals $t$: for a partition of mesh $\to 0$,
$\sum_j (W_{s_{j+1}} - W_{s_j})^2 \to t$ in mean square, since the sum
has mean $\sum_j (s_{j+1} - s_j) = t$ and variance
$2\sum_j (s_{j+1}-s_j)^2 \to 0$. Keeping the surviving terms:
\begin{equation*}
  \dd Y = \left(
    \frac{\partial \phi}{\partial t}
    + f\,\frac{\partial \phi}{\partial x}
    + \frac{g^2}{2}\,\frac{\partial^2 \phi}{\partial x^2}
  \right) \dd t
  \;+\; g\,\frac{\partial \phi}{\partial x}\; \dd W .
\end{equation*}
The half-Laplacian term, absent from the classical chain rule, is the
engine of everything that follows: it is why densities spread, why the
Fokker--Planck equation has a second-order term, and why exponential
transformations of SDEs pick up variance corrections. \hfill$\square$
\end{toolbox}

\section{The Fokker--Planck equation: transporting densities}
\label{sec:fokker-planck}

The SDE \eqref{eq:generic-sde} moves \emph{samples}; the corresponding
motion of the \emph{density} $p_t(x)$ is governed by the Fokker--Planck
(or forward Kolmogorov) equation. This change of viewpoint, from
trajectories to densities, is the same one Liouville's theorem performed
for Hamiltonian flow, and the derivation is parallel, with the It\^o
correction supplying a new, diffusive term.

\begin{toolbox}{Derivation of the Fokker--Planck equation}
Let $\phi$ be an arbitrary smooth test function of compact support, and
compute $\frac{\dd}{\dd t}\E[\phi(X_t)]$ in two ways.

\emph{Way 1: through the samples.} By It\^o's lemma,
\begin{equation*}
  \dd \phi(X_t)
  = \Big( f\,\phi' + \tfrac{g^2}{2}\,\phi'' \Big)\dd t
  + g\,\phi'\,\dd W_t ,
\end{equation*}
and taking expectations kills the martingale term:
\begin{equation*}
  \frac{\dd}{\dd t}\,\E[\phi(X_t)]
  = \E\Big[ f(X_t,t)\,\phi'(X_t) \Big]
  + \frac{g(t)^2}{2}\, \E\big[ \phi''(X_t) \big] .
\end{equation*}

\emph{Way 2: through the density.}
$\E[\phi(X_t)] = \int \phi(x)\, p_t(x)\,\dd x$, so
\begin{equation*}
  \frac{\dd}{\dd t}\,\E[\phi(X_t)]
  = \int \phi(x)\, \frac{\partial p_t(x)}{\partial t}\, \dd x .
\end{equation*}

Equate the two, and integrate Way 1 by parts to move all derivatives off
$\phi$ (boundary terms vanish by compact support): once for the drift
term, twice for the diffusion term,
\begin{equation*}
  \int \phi\, \frac{\partial p_t}{\partial t}
  = \int \phi \left[
      -\frac{\partial}{\partial x}\big( f\, p_t \big)
      + \frac{g^2}{2}\, \frac{\partial^2 p_t}{\partial x^2}
    \right] \dd x .
\end{equation*}
Since $\phi$ is arbitrary, the integrands must agree pointwise:
\begin{equation*}
  \boxed{\;
  \frac{\partial p_t}{\partial t}
  = -\frac{\partial}{\partial x}\big( f(x,t)\, p_t(x) \big)
  + \frac{g(t)^2}{2}\, \frac{\partial^2 p_t(x)}{\partial x^2} .
  \;}
\end{equation*}
The first term transports probability along the drift (as in Liouville);
the second spreads it diffusively. The diffusive term is the macroscopic
footprint of the It\^o correction. \hfill$\square$
\end{toolbox}

In vector form, and for later reference, the Fokker--Planck equation of
\eqref{eq:generic-sde} reads
\begin{equation}
  \partial_t p_t = -\nabla \cdot \big( f\, p_t \big)
  + \tfrac{1}{2} g^2\, \nabla^2 p_t .
  \label{eq:fokker-planck}
\end{equation}

\subsection{Stationary law of Langevin dynamics: closing the loop}

Apply the Fokker--Planck equation to the \emph{Langevin SDE} associated
with a potential $U$ (the continuous-time limit of the algorithm
\eqref{eq:ula}):
\begin{equation}
  \dd X_t = -\nabla U(X_t)\,\dd t + \sqrt{2\beta^{-1}}\;\dd W_t .
  \label{eq:langevin-sde}
\end{equation}
A stationary density solves
$0 = \nabla \cdot ( \nabla U\, p + \beta^{-1} \nabla p )$, i.e.\ the
probability current $J = -(\nabla U)\,p - \beta^{-1}\nabla p$ is
divergence-free; imposing the natural (decaying) boundary conditions
forces $J \equiv 0$, hence $\nabla p / p = -\beta \nabla U$, hence
\begin{equation}
  p_\infty(x) \;\propto\; e^{-\beta U(x)} :
  \label{eq:langevin-boltzmann}
\end{equation}
\emph{Langevin dynamics equilibrates to the Boltzmann distribution} of
Chapter~\ref{ch:statmech}. This is the two-line proof of the claim left
pending in Section~\ref{sec:mcmc}, and it explains the deep kinship of
sampling and physics: a Langevin sampler \emph{is} an overdamped physical
system at temperature $1/\beta$. Note also that the drift of
\eqref{eq:langevin-sde} is (up to $\beta$) the \emph{score} of the
target, $-\nabla U = \nabla \log p_\infty$: dynamics that sample a
distribution are driven by its score.

For the OU process, $U(x) = x^2/2$ and $\beta = 1$: the stationary law is
$\Nd(0,1)$, confirming the $t \to \infty$ limit of Section~\ref{sec:ou}.

